% Created 2023-12-28 jue 18:32
% Intended LaTeX compiler: pdflatex
\documentclass[11pt]{article}
\usepackage[utf8]{inputenc}
\usepackage[T1]{fontenc}
\usepackage{graphicx}
\usepackage{longtable}
\usepackage{wrapfig}
\usepackage{rotating}
\usepackage[normalem]{ulem}
\usepackage{amsmath}
\usepackage{amssymb}
\usepackage{capt-of}
\usepackage{hyperref}
\usepackage{minted}

\usepackage{parskip}
\usepackage[utf8]{inputenc} %% For unicode chars
\usepackage{placeins}
\usepackage[margin=2.50cm]{geometry}
\usepackage[T1]{fontenc}
\usepackage{mathpazo}
\linespread{1.05}
\usepackage[scaled]{helvet}
\usepackage{courier}
\hypersetup{colorlinks=true,linkcolor=blue}
\RequirePackage{fancyvrb}
\DefineVerbatimEnvironment{verbatim}{Verbatim}{fontsize=\small,formatcom = {\color[rgb]{0.5,0,0}}}
\usepackage{mdframed}
\BeforeBeginEnvironment{minted}{\begin{mdframed}}
\AfterEndEnvironment{minted}{\end{mdframed}}
\setcounter{secnumdepth}{3}
\date{}
\title{PHP-BasicExercises}
\hypersetup{
 pdfauthor={},
 pdftitle={PHP-BasicExercises},
 pdfkeywords={},
 pdfsubject={},
 pdfcreator={Emacs 28.1 (Org mode 9.6.12)}, 
 pdflang={Spanish}}
\begin{document}

\maketitle
\setcounter{tocdepth}{3}
\tableofcontents

\setlength\parindent{10pt}

\section{Practice PHP with these exercises}
\label{sec:orgfa69c1d}
Practicing with exercises is a great way to solidify your understanding
of PHP concepts. Here are some exercises that cover the topics from the
chapters on basics, conditionals, and loops. Each exercise is designed
to build upon the concepts you've learned.

\subsubsection{Exercise 1: Basic PHP Script}
\label{sec:org6629d79}
\textbf{Objective}: Write a simple PHP script that outputs "Hello, World!".

\subsubsection{Exercise 2: Variable Manipulation}
\label{sec:orgefb47da}
\textbf{Objective}: Create three variables: a string, an integer, and a
boolean. Perform the following tasks: - Concatenate the string with
another string. - Increment the integer by 10. - Toggle the boolean
value.

\subsubsection{Exercise 3: Conditional Greetings}
\label{sec:orga7a6ad6}
\textbf{Objective}: Write a PHP script that uses an if-else statement to output
"Good morning" if the time (in hours) is less than 12, and "Good
afternoon" otherwise.

\subsubsection{Exercise 4: Age Group Classifier}
\label{sec:org3cfee83}
\textbf{Objective}: Use an if-elseif-else statement to classify and echo an age
variable into categories (child, teenager, adult, senior). Define the
age categories as you see fit.

\subsubsection{Exercise 5: Loop through Numbers}
\label{sec:org3fe3d06}
\textbf{Objective}: Write a for loop that prints numbers from 1 to 50. For each
multiple of 3, print "Fizz" instead of the number, and for each multiple
of 5, print "Buzz". For numbers which are multiples of both 3 and 5,
print "FizzBuzz".

\subsubsection{Exercise 6: Array Iteration}
\label{sec:org1bd7c2e}
\textbf{Objective}: Given an array of names (\texttt{["Alice", "Bob", "Charlie"]}),
use a foreach loop to print each name with a greeting, like "Hello,
Alice!".

\subsubsection{Exercise 7: Factorial Calculator}
\label{sec:org5136aa4}
\textbf{Objective}: Write a PHP function to calculate and return the factorial
of a number using a for loop. Test the function with different numbers.

\subsubsection{Exercise 8: Prime Number Checker}
\label{sec:orga6de3c7}
\textbf{Objective}: Create a PHP script using a while loop that checks and
prints whether each number from 1 to 100 is a prime number.

\subsubsection{Exercise 9: Do-While Username Prompt}
\label{sec:org3e30a1e}
\textbf{Objective}: Simulate a do-while loop that continues to ask for a
username until the input is "admin".

\subsubsection{Exercise 10: Switch Case with Colors}
\label{sec:orge242a1e}
\textbf{Objective}: Use a switch case to handle different colors input by the
user. If the user inputs "red", "blue", or "green", output a message
like "You chose red." Include a default case that says "Unknown color".

\noindent\rule{\textwidth}{0.5pt}

These exercises are designed to challenge and reinforce your
understanding of PHP fundamentals, conditionals, and loops. Feel free to
adjust the difficulty or add new requirements to the exercises to suit
your learning needs!
\end{document}
